\section{The two laws on exact mesh intersections}\label{sec:methods}

% Describes core/src/gonthier.cpp, core/src/distributed.cpp and core/src/pcm.cpp at the commit of
% the battery. The constructions that this section only names are in Sections S1 and S2 of the
% supplement; if the laws change, change both.

Both laws are the Winkler foundation applied to the overlap $\region = A\cap B$ of two
undeformed bodies given by closed triangle meshes. The overlap is read as a bundle of columns
along a normal $\n$; a column of length $\ell$ stands for a spring compressed by $\ell$ and
presses with $K\ell$ per unit of area, with $K$ in N/m$^3$. A rate term of Hunt--Crossley form
multiplies the pressure by $1 - h\,v_n$, with $v_n = \n\cdot\vv$ the rate at which the bodies
separate and $h$ [s/m] the damping factor, and friction is Coulomb's, with coefficient $\mu$.
For a flat contact the normal force is $K V(1 - h\,v_n)$, with $V$ the volume of the overlap.
The parameters $K$, $h$ and $\mu$ belong to the pair of bodies. The laws differ in what they do
with the foundation: Gonthier's law lumps it into one force at the centroid of the overlap and
torques from its inertia tensor; the distributed law integrates it over the faces of the
overlap, so that the centre of pressure and the friction torque follow from the distribution.
Table~\ref{tab:notation} lists the notation; $\vv(\x)$ and $\w$ are the velocity field and the
angular velocity of $A$ relative to $B$.

\begin{table}[t]
\caption{Notation. Vectors are bold}\label{tab:notation}
\footnotesize
\begin{tabular}{@{}lp{0.36\linewidth}lp{0.36\linewidth}@{}}
\toprule
Symbol & Meaning & Symbol & Meaning \\
\midrule
$\region$, $\partial\region$ & overlap $A\cap B$; its boundary &
$K$, $h$ & stiffness [N/m$^3$] and damping factor [s/m] of the foundation \\
$\partial\region_A$, $\partial\region_B$ & parts of $\partial\region$ on the surface of $A$ and of $B$ &
$\mu$, $\mu_s$ & kinetic and static coefficients of friction \\
$V$, $\vect{c}$, $\tens{J}$ & volume, centroid and inertia tensor (unit density, about $\vect{c}$) of $\region$ &
$\ell(\x)$, $\varphi(\x)$ & column length along $\n$; damping factor $1 - h\,\n\cdot\vv(\x)$ \\
$\n$, $\m$ & pair normal (pushes $A$ out of $B$); unit normal of a face &
$p$, $\Fn$ & pressure; normal force \\
$v_n$, $\vv_t$ & $\n\cdot\vv$ (positive when separating); tangential part of $\vv$ &
$\q$, $\rr$ & instantaneous centre of the slip in the plane of a triangle; $\x - \q$ \\
$\omega_n$, $\w_t$ & $\n\cdot\w$; tangential part of $\w$ &
$D$, $\vect{R}$, $\vect{U}$, $\tens{W}$ & polar moments of a triangle about $\q$ \\
$r_g$ & radius of gyration of $\region$ about $\n$ &
$d$, $\tau$, $v_s$ & pre-sliding displacement, stick time constant, stick--slip velocity \\
$A_c$, $A_T$ & area of a flat contact; of a triangle &
$\dt$, $w$ & time step; inverse mass seen by a contact \\
\botrule
\end{tabular}
\end{table}

\subsection{The overlap and its normal}\label{sec:geometry}

The boundary of the overlap is the part of the surface of $A$ inside $B$ and the part of the
surface of $B$ inside $A$, $\partial\region = \partial\region_A\cup\partial\region_B$. Each is
obtained by clipping the triangles of one mesh against a binary space partition of the other
\cite{Sutherland1974,Thibault1987}, after bounding-volume trees \cite{vandenBergen1997} have
selected the faces involved, as in the 2018 thesis \cite{Vilela2018}. Every clipped polygon
adds its term to the integrals of the monomials up to second degree over $\region$, which the
divergence theorem reduces to sums over the boundary \cite{Mirtich1996,Eberly2002}. The result
is the volume $V$, the centroid $\vect{c}$ and the inertia tensor $\tens{J}$ of the overlap,
exact for the polyhedra given up to round-off: there is no sampling and no resolution
parameter. The pair normal is the direction of the vector area of the part of the boundary that
belongs to $B$,
\begin{equation}
  \vect{S}_B = \int_{\partial\region_B} \m \,\dd A , \qquad \n = \frac{\vect{S}_B}{|\vect{S}_B|} ,
  \label{eq:normal}
\end{equation}
which depends only on the curve where the two surfaces meet: when that contour is planar, $\n$
is the normal of its plane. The choice is not new: according to Elandt et
al.~\cite{Elandt2019}, Hasegawa and Fujii \cite{Hasegawa2003} and Luo and Nahon \cite{Luo2006}
define the direction of the patch in the same way.

Both laws use one normal per contact region, and the distributed law requires convex shapes, so
a non-convex body is given as convex pieces that fill it exactly, without overlaps or gaps, and
one pair is evaluated for every piece of one body near a piece of the other. Cutting a body
creates internal faces (seams) that close the overlap of a piece but are not surface of the
body; they are tagged when the body is built and carry no pressure, end no column and do not
enter the normal of the distributed law, which makes its result independent of the cut
(Section~S1.4). An approximate convex decomposition of a general mesh \cite{Wei2022coacd} is
not covered. Where $\vect{S}_B$ cancels, as for a body enclosed in another, the pair keeps its
last normal (Section~S1.5). Coincident faces, the acceleration structures and the cost of a
pair are in Sections~S1.1, S1.2 and~S1.14.

\subsection{Gonthier's volumetric law}\label{sec:gonthier}

The law \law{gonthier} implements the volumetric model of Gonthier's thesis
\cite{Gonthier2007thesis} on the exact $V$, $\vect{c}$ and $\tens{J}$. The normal force acts
along $\n$ at the centroid, with a rolling-resistance torque that is the moment of the damped
pressure to first order in the depth of the overlap:
\begin{equation}
  \Fn = K V \left(1 - h\,v_n(\vect{c})\right), \qquad
  \vect{M}_r = -K h\, \tens{J}\, \w_t ,
  \label{eq:gonthier-normal}
\end{equation}
with $\Fn$ clamped at zero. For the friction the thesis defines the radius of gyration of the
overlap about the normal, $r_g^2 = \n\cdot\tens{J}\,\n/V$, and an average surface velocity,
$v_\mathrm{avg}^2 = |\vv_t|^2 + r_g^2\,\omega_n^2$ with $\vv_t$ taken at the centroid, and gives
a Coulomb force and a Coulomb spinning torque that share one friction budget (the Contensou
effect):
\begin{equation}
  \vect{f}_t = -\mu\,\Fn\,\frac{\vv_t}{v_\mathrm{avg}}, \qquad
  \vect{M}_s = -\mu\,\Fn\, r_g^2\,\frac{\omega_n}{v_\mathrm{avg}}\,\n .
  \label{eq:gonthier-friction}
\end{equation}
Pure sliding gives $\mu\Fn$ and no torque, pure spinning the torque $\mu\Fn r_g$ and no force,
and the spin needs no parameter of its own. The lumping has a known price: for a flat circular
patch of radius $R$ and uniform depth, $r_g = R/\sqrt{2}$, whereas the exact Coulomb integral
of a uniform pressure gives the arm $2R/3$, so the spinning torque of the model is
$3/(2\sqrt{2}) = 1.06$ times the exact one. This 6\,\% belongs to the published model and not
to its implementation. The stick--slip of the thesis and the four points in which the
implementation departs from it, none of which changes the sliding or the spinning law, are in
Section~S1.10. The variant \law{gonthier-thesis} keeps the friction of the second author's 2018
thesis \cite{Vilela2018} (a sliding force that ignores the spin, a viscous spinning torque and
a bristle anchored at the centroid of the overlap, which is not a material point) so that both
can be reported.

\subsection{The distributed law}\label{sec:distributed}

The law \law{distributed} evaluates the same foundation point by point on the contact surface,
which is $\partial\region$ without the seams. Through a point $\x$ of a face of $A$ the column
runs along $+\n$, into $A$ and out of $B$; through a point of a face of $B$ it runs along $-\n$
(Fig.~\ref{fig:columns}); it ends at the first face of the surface of either body, so its
length $\ell(\x)$ is the minimum of a few affine functions of $\x$. The pressure is
\begin{equation}
  p(\x) = \gamma\,K\,\ell(\x)\,\varphi(\x), \qquad \varphi(\x) = 1 - h\,\n\cdot\vv(\x) ,
  \label{eq:pressure}
\end{equation}
where $\gamma = \tfrac12$ when the two ends of the column lie on different bodies, so that each
end carries half and the law is symmetric in $A$ and $B$ ($\gamma = 1$ in an enclosure,
Section~S1.6). The pressure acts along the normal of the face it is on. The rate in $\varphi$
is measured along the pair normal because $\n\cdot\vv$ is the same all along a column and is
the rate at which the column shortens; sliding along the tangent plane of $\n$ and spinning
about $\n$ therefore damp nothing, on any face (Section~S1.7).

\begin{figure}[t]
\centering
% Section of an overlap with its columns and the pair normal.
\begin{tikzpicture}[scale=1.0, font=\footnotesize, >=Stealth]
  % Body B (ground with a step, so that the contact is not a single plane)
  \coordinate (b1) at (-4.2,0);
  \coordinate (b2) at (4.2,0);
  \fill[gray!18] (-4.2,-1.6) rectangle (4.2,0);
  \draw[thick] (b1) -- (b2);
  \node[anchor=south west] at (-4.1,-1.55) {$B$};
  % Body A: a wedge-like convex polygon pressed into B, tilted.
  \coordinate (a1) at (-2.6,0.55);
  \coordinate (a2) at (-0.4,-0.75);
  \coordinate (a3) at (2.9,0.35);
  \coordinate (a4) at (2.2,2.3);
  \coordinate (a5) at (-2.0,2.4);
  \fill[gray!8] (a1) -- (a2) -- (a3) -- (a4) -- (a5) -- cycle;
  \draw[thick] (a1) -- (a2) -- (a3) -- (a4) -- (a5) -- cycle;
  \node at (0.2,1.7) {$A$};
  % Intersections of A's lower edges with y = 0.
  \coordinate (p) at ($(a1)!{0.55/1.30}!(a2)$);   % on edge a1-a2 at y=0
  \coordinate (q) at ($(a2)!{0.75/1.10}!(a3)$);   % on edge a2-a3 at y=0
  % Overlap
  \fill[pattern=north east lines, pattern color=gray!70] (p) -- (a2) -- (q) -- cycle;
  \draw[very thick] (p) -- (q);                    % part of dV on B
  \draw[very thick, densely dashed] (p) -- (a2) -- (q);   % part of dV on A
  % Columns along n = +y: from A's faces up to B's face, and from B's face down to A's faces.
  \foreach \t in {0.2,0.4,0.6,0.8} {
    \coordinate (s) at ($(p)!\t!(a2)$);
    \draw[thin] (s) -- (s |- b1);
  }
  \foreach \t in {0.15,0.3,0.45,0.6,0.75,0.9} {
    \coordinate (s) at ($(a2)!\t!(q)$);
    \draw[thin] (s) -- (s |- b1);
  }
  % One column labelled.
  \coordinate (x) at ($(a2)!0.3!(q)$);
  \fill (x) circle (1.1pt) node[below, inner sep=2.5pt] {$\x$};
  \draw[thin] ($(x)!0.5!(x |- b1)$) -- (1.3,-1.0) node[right, inner sep=1.5pt] {$\ell(\x)$};
  % Pair normal.
  \draw[->, very thick] (-3.4,0) -- ++(0,1.0) node[left] {$\n$};
  % Half pressures on both ends of the column.
  \draw[->, thick] (x) -- ++(-0.22,0.66);
  \draw[->, thick] ($(x |- b1)+(0,0.02)$) -- ++(0,0.7) node[above, inner sep=1.5pt] {$\tfrac12 K\ell$ each};
  % Labels
  \node[anchor=west] at (2.3,-0.35) {$\partial\region_B$ (face of $B$ inside $A$)};
  \draw[thin] (2.3,-0.35) -- ($(p)!0.93!(q)$);
  \node[anchor=east] at (-1.75,-0.95) {$\partial\region_A$ (faces of $A$ inside $B$)};
  \draw[thin] (-1.75,-0.95) -- ($(p)!0.75!(a2)$);
  \node at (-0.62,-0.2) {$\region$};
\end{tikzpicture}

\caption{Section of the overlap $\region$ of a body $A$ pressed into a body $B$. The pair
normal $\n$ is the direction of the vector area of $\partial\region_B$ (thick line). Columns run
along $\n$ from the faces of $A$ and along $-\n$ from the faces of $B$; the length $\ell(\x)$ of a
column is the distance to the first face of either body. Each end of a column carries half of
$K\ell$, as a pressure along the normal of its own face}
\label{fig:columns}
\end{figure}

Three constructions make the normal wrench exact for the polyhedra, with no quadrature error
and no dependence on the triangulation; one triangle per flat face is enough. Each polygon of
the contact surface is split into the convex regions where one face ends the column, on which
$\ell$ is affine (Section~S1.6). Where the affine factor $\varphi$ changes sign the region is
cut along $\varphi = 0$ and the tensile side carries no pressure, since there is no adhesion.
On what remains the pressure is the product of two affine functions, whose force and moment
over a triangle have a closed form; each triangle stores the linear pressure $\tilde p$ with
the same force and moment (Section~S1.7), which is the exact pressure whenever $\ell$ or
$\varphi$ is constant on it. Hasegawa and Sato \cite{Hasegawa2004} obtain the exactness of the
elastic term in another way, with a depth that is linear on the triangles of the intersection
polyhedron.

For a flat contact the law reduces to Gonthier's normal wrench: since $\n\cdot\vv$ is constant
along each column and affine in position,
\begin{equation}
  \int K\,\ell\,(1 - h\,\n\cdot\vv)\,\dd A
  = K\!\int_{\region} (1 - h\,\n\cdot\vv)\,\dd V
  = K V \left(1 - h\,\n\cdot\vv(\vect{c})\right),
  \label{eq:reduction}
\end{equation}
and the first moment of the damped pressure is, up to a term of the order of the depth, the
rolling resistance $\vect{M}_r$ of Eq.~\eqref{eq:gonthier-normal}: under a rolling body the
leading half of the patch is loaded more than the trailing half. The rolling resistance is a
property of the foundation with damping, not an addition to it, and the results have to account
for it in the problems whose rigid-body theory has none (Section~\ref{sec:classA}).

\subsection{Coulomb friction in closed form}\label{sec:friction}

Every triangle of the contact surface carries a linear pressure $p$ and slips with the relative
velocity field of the two bodies. The friction traction on $A$ is $-\mu\,p\,\vv_t/|\vv_t|$, with
$\vv_t$ the projection of $\vv$ on the plane of the triangle, and its integral is not a
polynomial: the direction of the slip changes over the triangle. It can nevertheless be done in
closed form, because the slip field of a rigid motion in a plane is a rotation about a point.
Let $\m$ be the unit normal of the triangle, $\x_0$ one of its vertices, $\omega_m = \w\cdot\m$
and $\vv_{t0}$ the tangential velocity at $\x_0$. Then
\begin{equation}
  \vv_t(\x) = \omega_m\,\m\times\rr ,
  \qquad \rr = \x - \q, \qquad \q = \x_0 + \frac{\m\times\vv_{t0}}{\omega_m} ,
  \label{eq:centre}
\end{equation}
where $\q$ is the instantaneous centre of the slip in the plane of the triangle; it need not
lie inside the triangle (Fig.~\ref{fig:slip}). Writing the pressure about the centre,
$p(\x) = p_q + \nabla p\cdot\rr$, with $\sigma = \operatorname{sgn}\omega_m$, the friction force
on $A$ and its moment about $\q$ are
\begin{align}
  \vect{F}_f &= -\mu\,\sigma\,\m\times\int p\,\frac{\rr}{|\rr|}\,\dd A
     = -\mu\,\sigma\,\m\times\left(p_q\,\vect{U} + \tens{W}\,\nabla p\right), \label{eq:friction-force}\\
  \vect{M}_f &= -\mu\,\sigma\int p\,\rr\times\frac{\m\times\rr}{|\rr|}\,\dd A
     = -\mu\,\sigma\left(p_q\,D + \nabla p\cdot\vect{R}\right)\m , \label{eq:friction-moment}
\end{align}
with four \emph{polar moments} of the triangle $T$ about $\q$,
\begin{equation}
  D = \int_T |\rr|\,\dd A, \quad
  \vect{R} = \int_T \rr\,|\rr|\,\dd A, \quad
  \vect{U} = \int_T \frac{\rr}{|\rr|}\,\dd A, \quad
  \tens{W} = \int_T \frac{\rr\,\rr^{\mathsf T}}{|\rr|}\,\dd A .
  \label{eq:moments}
\end{equation}
$D$ and $\vect{U}$ are the moments that Jia \cite{Jia1997} gives for a polygon under uniform
pressure; $\vect{R}$ and $\tens{W}$ are what the linear term of the pressure adds.

\begin{figure}[t]
\centering
% Slip field about the instantaneous centre, and the polar decomposition for one edge.
\begin{tikzpicture}[scale=1.0, font=\footnotesize, >=Stealth]
  % ---- Left: slip field -------------------------------------------------------
  \begin{scope}
    \coordinate (q) at (0,0);
    \coordinate (p0) at (1.3,0.5);
    \coordinate (p1) at (4.3,0.2);
    \coordinate (p2) at (2.9,2.9);
    \fill[gray!12] (p0) -- (p1) -- (p2) -- cycle;
    \draw[thick] (p0) -- (p1) -- (p2) -- cycle;
    \fill (q) circle (1.3pt) node[below left] {$\q$};
    % Arcs of the rotation.
    \foreach \rad in {2.0,3.0,3.8} {
      \draw[gray, thin, densely dotted] (5:\rad) arc (5:60:\rad);
    }
    % Slip vectors at sample points: perpendicular to r, length proportional to |r|.
    \foreach \px/\py in {1.9/0.7, 3.2/0.6, 2.4/1.5, 3.5/1.3, 2.9/2.2, 2.0/1.0} {
      \draw[->] (\px,\py) -- ++({-0.22*\py},{0.22*\px});
    }
    \coordinate (x) at (3.5,1.3);
    \draw[thin] (q) -- node[below, pos=0.55, inner sep=2pt] {$\rr$} (x);
    \fill (x) circle (1pt) node[right, inner sep=2pt] {$\x$};
    \node[anchor=east] at (2.2,2.6) {$\vv_t = \omega_m\,\m\times\rr$};
    \draw[->] (0.55,0) arc (0:70:0.55) node[above left, inner sep=1pt] {$\omega_m$};
  \end{scope}
  % ---- Right: polar coordinates for one edge ----------------------------------
  \begin{scope}[xshift=6.6cm]
    \coordinate (q) at (0,0.2);
    \coordinate (a) at (3.0,-0.5);
    \coordinate (b) at (3.0,2.9);
    \coordinate (f) at (3.0,0.2);
    \fill[gray!12] (q) -- (a) -- (b) -- cycle;
    \draw[thick] (a) -- (b);
    \draw (q) -- (a) node[below] {$\vect{a}$};
    \draw (q) -- (b) node[above] {$\vect{b}$};
    \draw[dashed] (q) -- node[below, pos=0.45, inner sep=2pt] {$\eta$} (f);
    \draw (2.8,0.2) -- (2.8,0.4) -- (3.0,0.4);
    \fill (q) circle (1.3pt) node[left] {$\q$};
    \draw[->, thick] (q) -- ++(0.9,0) node[below, inner sep=2pt] {$\vect{e}_1$};
    \draw[->, thick] (q) -- ++(0,0.9) node[left, inner sep=2pt] {$\vect{e}_2$};
    % A ray at angle theta.
    \coordinate (r) at (3.0,1.7);
    \draw[thin] (q) -- (r);
    \draw[->] ($(q)+(1.5,0)$) arc (0:26.5:1.5);
    \node at (1.75,0.6) {$\theta$};
    \node[anchor=west] at (3.05,1.7) {$r = \eta\sec\theta$};
    \node[anchor=west] at (3.05,-0.2) {$\theta_a$};
    \node[anchor=west] at (3.05,2.6) {$\theta_b$};
  \end{scope}
\end{tikzpicture}

\caption{Left: the slip field on a triangle is a rotation about the instantaneous centre $\q$,
with speed $|\omega_m||\rr|$ perpendicular to $\rr$. Right: the triangle is the signed sum of
the three triangles formed by $\q$ and each edge; on the triangle $(\q, \vect{a}, \vect{b})$,
polar coordinates about $\q$ with the angle $\theta$ measured from the perpendicular $\vect{e}_1$
to the edge give $r \le \eta\sec\theta$, with $\eta$ the distance from $\q$ to the line of the edge}
\label{fig:slip}
\end{figure}

The moments are evaluated edge by edge: the triangle is the signed sum of the three triangles
formed by $\q$ and each edge, and in polar coordinates about $\q$ each of them leaves four
elementary integrals whose primitives are an inverse hyperbolic sine and powers of a square
root. Appendix~\ref{app:polar} gives the result, Eq.~\eqref{eq:edge}, and its derivation. The
sign of each contribution comes out of the order of the limits, so no case distinction is
needed for $\q$ inside or outside the triangle, and $\q$ on an edge or at a vertex is not a
singular case. The cost is two square roots and two inverse hyperbolic sines per edge. A
script distributed with the sources re-implements the formulas as printed and compares them
with quadrature that does not use them, for $\q$ outside, inside, on an edge and at a vertex;
they agree to $10^{-14}$ where the integrand is analytic.
% The number of this paragraph is the one paper/checks/check_polar_moments.py prints; it is not a battery result.

The friction is exact for the linear pressure $\tilde p$ of each flat triangle, whatever the
triangulation. It is approximate where $\tilde p$ is not the true pressure (both $\ell$ and
$\varphi$ vary over a triangle), and a uniform-slip formula replaces it when the triangle does
not rotate in its plane or the centre is far (Section~S1.9).

\subsection{Stiction}\label{sec:stiction}

Coulomb friction regularized by velocity alone lets a loaded body creep. Both laws therefore
hold a contact at rest with a bristle: a spring--damper on a tangential deflection $\vect{z}$,
blended with the sliding friction by a weight $s$,
\begin{equation}
\begin{gathered}
  \vect{f} = s\,\vect{f}_\mathrm{stick} + (1-s)\,\vect{f}_\mathrm{slide},
  \qquad \vect{f}_\mathrm{stick} = -k\,\vect{z} - c\,\vv_t, \\
  s = \exp\!\left(-u^2/v_s^2\right),
  \qquad k = \frac{\mu_s\,|\Fn|}{d}, \qquad c = k\,\tau ,
\end{gathered}
  \label{eq:stick}
\end{equation}
and the same for a torsional deflection about $\n$. The speed $u$ is Gonthier's
$v_\mathrm{avg}$ in \law{gonthier} and its counterpart for the pressure distribution in
\law{distributed}. The deflection is the slip of the material point of $A$ at the friction
point, accumulated between evaluations, so that rolling without slipping loads no bristle. The
stiffness is proportional to the load, as in the Dahl and LuGre models and in Gonthier's
bristle \cite{Gonthier2007thesis}: every contact reaches its static limit after the same
displacement $d = 0.1$\,mm, with $\tau = 5$\,ms and $v_s = 24.5$\,mm/s for every contact of
every scenario. The stick force is capped at the static limit, which in \law{distributed} is
the limit surface of the patch approximated by the ellipse of Howe and Cutkosky
\cite{Howe1996}, with the spinning capacity of the pressure distribution from the polar moments
of Eq.~\eqref{eq:friction-moment}, and in \law{gonthier}, as published, a box: force and torque
are limited separately (Section~S1.11).

\subsection{Explicit coupling with the engines and its step limit}\label{sec:coupling}

The contact model is a function: it receives the poses and velocities of the bodies and the
friction state of every pair and returns one wrench per body and the next state. The backends
for Drake 1.58, Project Chrono 10.0 and MuJoCo 4.0 switch the engine's own contact off, apply
the wrenches as external forces and advance (Section~S2.1). In the three the wrench is
\emph{explicit}: computed from the state at the start of the step and held during it, so for
contact the three steppers are the same semi-implicit Euler scheme. Running one law on three
engines therefore checks the implementation on code bases that share nothing, and cannot
separate the error of the law from that of the time discretization, which the step series does
(Section~\ref{sec:convergence}).

An explicit step has a stability limit. A spring $k$ and a damper $c$ acting on an inverse mass
$w$ are stable if and only if $k\,w\,\dt^2 + 2\,c\,w\,\dt < 4$ and $c\,w\,\dt < 2$
(Section~S2.2), and the contributions of the contacts of a body add. For the normal stiffness
of a flat contact of area $A_c$ this is
\begin{equation}
  K A_c\,\dt^2 \sum w < 4 ,
  \label{eq:stability}
\end{equation}
and $< 2$ with the dampers at their bounds: a body of mass $m$ held between two flat contacts,
for which $\sum w = 2/m$, needs $\dt < \sqrt{m/(K A_c)}$. The laws do not change $K$, so
Eq.~\eqref{eq:stability} is the limit of the coupling. The terms proportional to the load (the
stick spring and damper, the normal damper, the slope of the friction across the stick--slip
transition) would impose limits of their own under loads of several times a body's weight; the
laws bound them per unit of load, which changes nothing below a load threshold and, above it,
applies a softer stick, a wider transition and less damping than the parameters say, as
functions of the step (Section~S2.3). An implicit coupling needs the derivatives of the wrench,
which the laws do not provide yet.

\subsection{Baselines from the literature on the same geometry}\label{sec:baselines}

% Describes core/src/pcm.cpp and DistributedParams::friction_quadrature; formulations, sources
% and departures in docs/baselines.md.

Three further methods of the library, run like the two laws with the same $K$, $h$ and $\mu$,
show what the closed-form friction buys and how the closest published law behaves. Two variants
of \law{distributed} replace only the closed form by the quadratures of the literature: the
\emph{vertex rule} of Hasegawa and Sato \cite{Hasegawa2004}, which evaluates the friction at
the three vertices of each triangle and interpolates it linearly, and the \emph{one-point rule},
which applies $-\mu\int p\,\dd A$ along the slip direction at the pressure centroid of each
triangle, as the polygonal contact model, Simbody's elastic foundation and the discrete
hydroelastic model do per element \cite{Hippmann2024,Sherman2011,Masterjohn2022}. When the slip
direction is uniform over a triangle the three rules coincide, so they can differ only where a
contact spins about the normal of its faces. The law \law{pcm} is a per-face elastic foundation
after the polygonal contact model \cite{Hippmann2004,Hippmann2024} and Simbody
\cite{Sherman2011,SimbodyEFF}: every triangle of the surface of one body inside the other is a
spring with one sample of depth, velocity and slip, a normal force
$\tfrac12 K A_T\,x\,(1 + h\,\dot x)$ for a penetration $x$, and a regularized Coulomb friction
without stiction or state, as in its sources. Its elements are the triangles of the exact
contact surface and its depth is measured to the nearest point of the other surface; these and
its other departures from the publications are in Section~S1.13.
